\section{Preliminaries}\label{prelim}


In this subsection, we would like to first revisit the fact that
\begin{equation} \label{cupcap}
        MS_3(n; 2) \leq \binom{2n}{n} + 1
    \end{equation}
    holds for every $n \geq 4$. In \cite{ES35}, Erd\H{o}s and Szekeres originally proved this result in the language of cups and caps by establishing a recursive inequality in the spirit of the classical recursive inequality for diagonal Ramsey numbers (also originally from \cite{ES35}). In \cite{CK71} and \cite{MS14}, however, Chv\'atal--Koml\'os and Moshkovitz--Shapira note that it is also possible to prove \eqref{cupcap} by appropriately extending the Seidenberg argument \cite{Seid59}. Our proof of Theorem \ref{main3} (and that of Theorem \ref{maink}) will also adopt this perspective, so for the sake of self-containment we will now recall this beautiful idea below. 

   \subsection{Seidenberg's argument in uniformity three.} \label{Seidenberg} Let $\Delta: [N]^{(3)} \to \{\operatorname{red},\operatorname{blue}\}$ be a two-coloring of the triples of $[N]$ with no tight monotone monochromatic path of length $n$.
    For each $u < v$, define the vector
    \begin{equation*}
        P(u,v) = \begin{bmatrix}
            \l_{\text{red}}(u,v) + 1\\
            \l_{\text{blue}}(u,v) + 1
        \end{bmatrix},
    \end{equation*}
    where $\l_{\text{red}}(u,v)$ and $\l_{\text{blue}}(u,v)$ denote the lengths of the longest tight monotone paths in red and blue, respectively, which end at $u, v$.
    Since there is no tight monotone path of length $n$, note that the point $P(u,v)$ lies in the grid $[n]^2$. Now, consider the map
    \begin{equation*}
        \Theta: v \longmapsto D(v) = \{\vx \in [n]^2: \vx \leq P(u,v) \text{ for some } u < v\},
    \end{equation*}
    which assigns elements of $[N]$ down-sets of $[n]^2$. Here the $\leq$ relation on $[n]^2$ denotes the usual ``less than or equal to in both coordinates" partial order on $[n]^2$. One can check that $D(v)$ is indeed a down-set: for every $\vx \in D(v)$ we have that $\vx' \in D(v)$ for every $\vx' \leq \vx$. 
    
    The main observation is that the map $\Theta$ is injective: if $u < v$ and $D(u) = D(v)$, we have $P(u,v) \in D(v) = D(u)$, so by definition $P(u,v) \leq P(w,u)$ for some $w < u$; but this is a contradiction since the longest path ending at $u, v$ in color $\Delta(w,u,v)$ is strictly longer than the one ending at $w,u$.
    Since the number of down-sets in $[n]^2$ equals the number of antichains in $[n]^2$ which is $\binom{2n}{n}$, we have that $N \leq \binom{2n}{n}$. We refer to \cite{MS14} for more details regarding the latter claims. 

    %\bigskip

    The main idea behind Theorem \ref{main3} is that in order to upper bound $A_{3}(n;r,s)$ one can construct an injection of $[N]$ into a significantly smaller space than the collection of down-sets of $[n]^2$. To define this map, we will require the concept of a dominating set in a majority tournament. 
    
\subsection{Majority tournaments}

A {\it{tournament}} is a directed graph $T=(V,E)$ where every two vertices $u,v \in V$ are connected by one edge, i.e. $|E \cap \left\{(u,v),(v,u)\right\}|=1$. If the edge is $(u,v)$ (denoted by $u \to v$) we say that $u$ {\it{dominates}} $v$. 

\begin{definition}
    For an integer $K \geq 2$, a tournament $T$ is called a $K$-\textit{majority tournament} if there exist $2K-1$ linear (total) orders $\prec_1, ..., \prec_{2K-1}$ on the set of vertices of $T$ such that $u \to v$ in $T$ if and only if $v \prec_i u$ holds for at least $K$ values of $i \in \left\{1,\ldots,2K-1\right\}$. 
\end{definition}

These types of tournaments have been studied in social science~\cite{Mcg53} as well as combinatorics, where previous researchers studied extremal problems regarding the sizes of induced transitive subtournaments \cite{MSW11} and dominating sets \cite{ABKKW06}. If $T$ is a tournament and $X,Y$ are subsets of vertices of $T$, we say that $X$ {\it{dominates}} $Y$ if for every vertex $y$ in $Y \setminus X$ there exists an $x \in X$ such that $x$ dominates $y$. If $Y$ is the full set of vertices of $T$, then we call $X$ a {\it{dominating set}} of the tournament $T$. The {\it{domination number}} $f(T)$ of $T$ is the smallest cardinality of a dominating set of $T$. It is a classical fact (often attributed to Erd\H{o}s, see for example \cite{Moon68}) that every tournament $T$ on $n$ vertices satisfies $f(T) \leq \lceil \log_{2}{n} \rceil$. It is also known \cite{AS16} that random tournaments on $n$ vertices do not contain any smaller dominating sets, so this estimate cannot be improved upon in general. Nevertheless, it turns out that $K$-\textit{majority tournaments} behave quite differently than random tournaments when it comes to their domination number $f(T)$. In \cite{ABKKW06}, Alon, Brightwell, Kierstead, Kostochka, Winkler proved the following result:

\begin{theorem} \label{kmaj}
If $T$ is a $2$-majority tournament, then 
$$f(T) \leq 3.$$
In general, for every $K \geq 2$, if $T$ is a $K$-majority tournament, then
$$f(T) \leq C K \log K$$
holds for an absolute constant $C>0$. 
\end{theorem}

It turns out it is this phenomenon which is primarily responsible for the drop in the tower height in our estimates from Theorem \ref{main3} and Theorem \ref{maink} when compared to the estimate from Theorem \ref{MS}. 

%%%%%%%%%%%%%%%%%%%%%%%%%%%%%%%%%%%%%%%%%%%%%%%%%%%%%%%%%%%%%%%%%%%%%%%%%%%%%%%%%%%%%%%%%%%%%%%%%%%%%%%%%%%%%%%%%%%%%%%%%%%%%%%%%%%%%%%%%%%%%%%%%%%%%%%%%%%%%%%%%%%%%%%%%%%%%%%%%%%%%%%%%%%%%%%%%%%%%%%%%%%%%%%%%%%%%%%%

\section{Proof of Theorem~\ref{main3}}\label{s3}

\subsection{The proof for $r=3$, $s=2$.}

We first prove Theorem~\ref{main3} for $r=3$ and $s=2$, in order to showcase the main idea in the clearest way possible. Our proof for this regime actually gives the stronger bound $A_3(n;3,2) \leq n^9$.

    Let $\Delta: [N]^{(3)} \to [3]$ be a 3-coloring of the triples of $[N]$ in which every tight monotone path of length $n$ receives all three colors.
    Define for each $u < v$ in $[N]$ the vector
    \begin{equation*}
        P(u,v) = \begin{bmatrix}
            \l_{1,2}(u,v) + 1\\
            \l_{2,3}(u,v) + 1\\
            \l_{1,3}(u,v) + 1
        \end{bmatrix} \in [n]^3,
    \end{equation*}
    where $\l_{i,j}(u,v)$ is the length of the longest tight monotone path ending at $u, v$ which uses only colors $i$ and $j$.
    
    Like in the argument from Subsection \ref{Seidenberg}, let us consider the map
    \begin{eqnarray*}
        \Theta_1: v \longmapsto S(v) = \{P(u,v) \in [n]^3: u < v\},
    \end{eqnarray*}
    associating elements of $[N]$ with subsets of $[n]^3$. 
    To define an injection from $[N]$ into a space of size smaller than ${2n \choose n}$, we define a 2-majority tournament on the elements of $S(v)$ with linear orders $\prec_1, \prec_2, \prec_3$ as follows.
    If $\vx = (x_1,x_2,x_3), \vy = (y_1,y_2,y_3) \in S(v)$, set
    \begin{align*}
        & 1. \hspace{0.25cm} \Vec{x} \prec_1 \Vec{y} \text{ if } \begin{cases}
            x_1 < y_1; \text{ or}\\
            x_1=y_1, x_2 < y_2; \text{ or}\\
            x_1=y_1, x_2=y_2, x_3 < y_3
        \end{cases}\\
        & 2. \hspace{0.25cm} \Vec{x} \prec_2 \Vec{y} \text{ if } \begin{cases}
            x_2 < y_2; \text{ or}\\
            x_2=y_2, x_3 < y_3; \text{ or}\\
            x_2=y_2, x_3=y_3, x_1 < y_1
        \end{cases}\\
        & 3. \hspace{0.25cm} \Vec{x} \prec_3 \Vec{y} \text{ if } \begin{cases}
            x_3 < y_3; \text{ or}\\
            x_3=y_3, x_1 < y_1; \text{ or}\\
            x_3=y_3, x_1=y_1, x_2 < y_2
        \end{cases}
    \end{align*}
    
    In words, $\prec_i$ is the lexicographic order on $S(v) \subseteq [n]^3$ in which the $i$th coordinate is the most important, followed by the $(i+1)$-st and then the $(i+2)$-nd (with addition modulo 3).
    By Theorem \ref{kmaj}, $S(v)$ has a total dominating set $D(v)$ of size at most three, i.e., for every $\vy \in S(v)$, there exists $\vx \in D(v)$ with $\vx \to \vy$ in the tournament defined above. 
    
    Now consider a second map
    \begin{equation*}
        \Theta_2: v \longmapsto D(v) \subseteq S(v).
    \end{equation*}
    Since $D(v) \subseteq S(v) \subseteq [n]^3$, note that $\Theta_{2}$ is a map from $[N]$ to the set of subsets of size at most $3$ of $[n]^3$. We claim that $\Theta_2$ is also an injection.
    Suppose, on the contrary, that $u < v$ but $D(u) = D(v)$.
    Then each element $\Vec{x} \in D(v) = D(u) \subseteq S(u)$ by definition has the form $\Vec{x} = P(w,u)$ for some $w < u$;
    but $\Vec{x} = P(w,u) \prec_i P(u,v)$ for at least the two coordinates $i$ corresponding to paths in which the color $\Delta(w,u,v)$ may be used, and so $P(u,v) \to \vx$ for all $\vx \in D(v)$, contradicting the fact that $D(v)$ totally dominates $S(v)$ in the 2-majority tournament defined by the $\prec_i$'s.

    Therefore, 
    %since we may take each $D(v)$ to have size at most three, 
    we have $N \leq \binom{n^3}{3}+\binom{n^3}{2} +\binom{n^3}{1}+\binom{n^3}{0} \le n^9$ as promised. \qed
    %\begin{flushright}
    %$\square$
    %\end{flushright}

%\vspace{-10mm}

\subsection{General $r$ and $s$.}

Now we prove Theorem~\ref{main3} for all combinations $r, s$ with $s > r/2$.
The proof is a straightforward generalization of the one given in the previous subsection.

    Let $\Delta: [N]^{(3)} \to [r]$ be an $r$-coloring of the triples of $[N]$ in which every tight monotone path on $n$ edges receives at least $s+1$ colors (i.e., no such path uses at most $s$ colors).
    As before, define for each $u < v$ the vector
    \begin{equation*}
        P(u,v) = \langle 1 + \l_S(u,v)\rangle_{S \in [r]^{(s)}} \in [n]^{\binom{r}{s}},
    \end{equation*}
    where $\l_S(u,v)$ denotes the length of the longest monotone path ending at $u$ and $v$ which only uses colors from the set $S \in [r]^{(s)}$,
    and consider the map
    \begin{equation*}
        \Theta_1: v \longmapsto S(v) = \{P(u,v) \in [n]^{\binom{r}{s}}: u < v\}
    \end{equation*}
    associating elements of $[N]$ with subsets of $[n]^{\binom{r}{s}}$.

    Let $K$ be the positive integer satisfying $\binom{r}{s} = 2K-1$ if $\binom{r}{s}$ is odd and $\binom{r}{s} = 2(K-1)$ if $\binom{r}{s}$ is even.
    For each $v \in [N]$, we define a $K$-majority tournament on the elements of $S(v)$ with linear orders $\prec_1, \prec_2, \dots, \prec_{\binom{r}{s}}$ as follows.
    If $\vx = (x_1,x_2, \dots, x_{\binom{r}{s}}), \vy = (y_1,y_2, \dots, y_{\binom{r}{s}}) \in S(v)$ and $1 \leq i \leq \binom{r}{s}$ (with addition in the indices considered modulo $\binom{r}{s}$), set  
    \begin{equation*}
        \vx \prec_i \vy \text{ if } \begin{cases}
            x_i < y_i; \text{ or}\\
            x_i=y_i, x_{i+1} < y_{i+1}; \text{ or}\\
            x_i=y_i, x_{i+1} = y_{i+1}, x_{i+2} < y_{i+2}; \text{ or}\\
            \vdots\\
            x_i=y_i, x_{i+1} = y_{i+1}, x_{i+2} = y_{i+2}, \dots, x_{i+\binom{r}{s}-2} = y_{i+\binom{r}{s}-2}, x_{i+\binom{r}{s}-1} < y_{i+\binom{r}{s}-1}
        \end{cases}
    \end{equation*}
    Note that if $\binom{r}{s} = 2K-1$ is odd then this defines a $K$-majority tournament on $S(v)$.
    If $\binom{r}{s} = 2(K-1)$ is even, then we will force the above construction to be a $K$-majority tournament by adding an additional arbitrary linear order $\prec_{*}$ with no relation to the others.
    \begin{comment}
    Note that by our assumption that $s > r/2$,
    \begin{align*}
        \binom{r}{s} = \frac{r}{s} \cdot \binom{r-1}{s-1} < 2 \cdot \binom{r-1}{s-1},
    \end{align*}
    so these orders do in fact define an $\binom{r-1}{s-1}$-majority tournament on $S(v)$.
    \end{comment}
    By Theorem \ref{kmaj}, each $S(v)$ has a total dominating set $D(v)$ in the tournament defined above, i.e., for every $\vy \in S(v)$, there exists $\vx \in D(v)$ with $\vx \to \vy$, with $|D(v)| \le C K \log K$.

    
    Now consider a second map
    \begin{equation*}
        \Theta_2: v \longmapsto D(v) \subseteq S(v).
    \end{equation*}
    We argue that $\Theta_2$ is again an injection.
    To see this, assume otherwise that for some $u < v$ in $[N]$, $D(u) = D(v)$.
    %Observe that for all $w < u$, 
    %We know that there is some $\vy \in S(v)$ such that for every $\vx \in T(u) \subseteq D(u)$, $\vx \prec_{\binom{r-1}{s}} y$; but $T(v) = T(u)$, so this holds for every $\vx \in T(v)$ as well.
    Each $\vx \in D(v) = D(u) \subseteq S(u)$ has the form $P(w,u)$ for some $w < u$.
    Now $\vx = P(w,u) \prec_i P(u,v)$ in all of the $\binom{r-1}{s-1}$ coordinates $i \in [\binom{r}{s}]$ which correspond to subsets $S \subseteq [r]$ containing the color $\Delta(w,u,v)$.
    By our hypothesis that $s > r/2$, we have
    \begin{equation*}
        2(K-1) \leq \binom{r}{s} = \frac{r}{s} \cdot \binom{r-1}{s-1} < 2 \cdot \binom{r-1}{s-1},
    \end{equation*}
    which gives (regardless of parity, since all values are integers) that 
    \begin{equation*}
        \binom{r-1}{s-1} \geq K.
    \end{equation*}
    Thus, each such $\vx$ has $\vx \prec_i P(u,v)$ for at least $K$ of the orders $\prec_i$, $1 \leq i \leq \binom{r}{s}$, and so $P(u,v) \to \vx$ in the $K$-majority tournament on $S(v)$ for all $\vx \in D(v)$, contradicting the fact that $D(v)$ is a total dominating set of $S(v)$.

    Therefore, since each $D(v)$ has size at most $C K \log K = C' \binom{r}{s} \log \binom{r}{s}$, we have
    \begin{equation*}
        N \leq \binom{n^{\binom{r}{s}}}{\le C' \binom{r}{s} \log \binom{r}{s}} \leq n^{C' \binom{r}{s}^2 \log \binom{r}{s}}.
    \end{equation*}
    \qed
%\begin{flushright}
%$\square$
%\end{flushright}
%%%%%%%%%%%%%%%%%%%%%%%%%%%%%%%%%%%%%%%%%%%%%%%%%%%%%%%%%%%%%%%%%%%%%%%%%%%%%%%%%%%%%%%%%%%%%%%%%%%%%%%%%%%%%%%%%%%%%%%%%%%%%%%%%%%%%%%%%%%%%%%%%%%%%%%%%%%%%%%%%%%%%%%%%%%%%%%%%%%%%%%%%%%%%%%%%%%%%%%%%%%%%%%%%%%%%%%%

\vspace{-10mm}

\section{Proof of Theorem~\ref{maink}}\label{sk}

The argument is an extension of the one in uniformity three, but we no longer define an injection from the singletons into the subsets of $[n]^{\binom{r}{s}}$ of size at most $C \binom{r-1}{s-1} \log \binom{r-1}{s-1}$.
Instead, we define such a map on $[N]^{(k-2)}$ and show that it contains no 2-edge monotone path in which both edges are mapped to the same set, or, in other words, we define a coloring of $[N]^{(k-2)}$ with $n^{C \binom{r}{s} \binom{r-1}{s-1} \log \binom{r-1}{s-1}}$ colors and no monochromatic monotone path of length two.
This allows us to use Theorem~\ref{MS} to get an improved upper bound for $A_k(n;r,s)$ over the trivial $MS_k(n;\lceil r/s \rceil)$.
Similar ideas were used in \cite{FPSS12, MS14} to upper bound $MS_k(n,r)$.\\

\noindent \textit{Proof of Theorem~\ref{maink}}.
    Let $\Delta: [N]^{(k)} \to [r]$ be an $r$-coloring of the $k$-subsets of $[N]$ in which every tight monotone path on $n$ edges receives at least $s+1$ colors.
    For $v_1 < v_2 < \dots < v_{k-1}$, define
    \begin{equation*}
        P(v_1,v_2,...,v_{k-1}) =
            \langle 1 + \l_S(v_1, v_2, ...,v_{k-1})\rangle_{S \in [n]^{(r)}} \in [n]^{\binom{r}{s}},
    \end{equation*}
    where as above $\l_S(v_1,...,v_{k-1})$ denotes the length of the longest tight monotone path ending at $v_1, ..., v_{k-1}$ which only uses colors in the set $S \in [r]^{(s)}$.
    
    Consider the function
    \begin{align*}
         \Theta_1: \{v_1,...,v_{k-2}\} \longmapsto S(v_1,...,v_{k-2}) = \{P(u,v_1,...,v_{k-2}) \in [n]^{\binom{r}{s}}: u < v_1\}
    \end{align*}
    mapping elements of $[N]^{(k-2)}$ into subsets of $[n]^{\binom{r}{s}}$.
    We can define the exact same $K$-majority tournament as above on each $S(v_1,...,v_{k-2})$ and obtain a total dominating set $D(v_1,...,v_{k-2}) \subseteq S(v_1,...,v_{k-2})$ of size at most $C K \log K$.
    With this, again consider a second map
    \begin{equation*}
        \Theta_2: \{v_1,...,v_{k-2}\} \longmapsto D(v_1,...,v_{k-2}) \subseteq S(v_1,...,v_{k-2}).
    \end{equation*}
    Now let $S \subseteq [n]^{\binom{r}{s}}$ have size $|S| \le C K \log K$.
    We claim that $\Theta_2^{-1}(S)$ has no tight $(k-2)$-uniform monotone path of length two.
    Indeed, if it had such a path on vertices $v_1 < v_2 < ... < v_{k-1}$, then each $\Vec{x} \in D(v_2,...,v_{k-1}) = S = D(v_1,...,v_{k-2}) \subseteq S(v_1,...,v_{k-2})$ by definition has $\Vec{x} = P(v_0,v_1,...,v_{k-2})$ for some $v_0 < v_1$; but then
    \begin{equation*}
        \Vec{x} = P(v_0,...,v_{k-2}) \prec_i P(v_1,...,v_{k-1}) \in S(v_2,...,v_{k-1})
    \end{equation*}
    in the $\binom{r-1}{s-1} \geq K$ coordinates $i$ corresponding to paths which are allowed to contain the color $\Delta(v_0,...,v_{k-1})$,
    and so $P(v_1, \dots, v_{k-1}) \to \vx$ for all $\vx \in D(v_2, \dots, v_{k-1})$, contradicting the fact that $D(v_2,...,v_{k-1})$ was chosen as a total dominating set for $S(v_2,...,v_{k-1})$.
    
    In other words, then, the partition
    \begin{equation*}
        [N]^{(k-2)} = \bigcup_{S \subseteq [n]^{\binom{r}{s}}, |S| \le C K \log K} \Theta_2^{-1}(S)
    \end{equation*}
    defines an $\binom{n^{\binom{r}{s}}}{\le C K \log K}$-coloring of the $(k-2)$-tuples of $[N]$ with no monochromatic monotone path of length two.
    Thus, by Theorem \ref{MS}, when $k \geq 4$ we have
    \begin{equation*}
        N 
        \leq MS_{k-2} \left( 2; n^{C' \binom{r}{s}^2 \log \binom{r}{s}} \right) 
        \leq T_{k-4}\left(2^{n^{C' \binom{r}{s}^2 \log \binom{r}{s}}}\right)
        = T_{k-3} \left( n^{C' \binom{r}{s}^2 \log \binom{r}{s}} \right).
    \end{equation*}
\qed
